%%%%%%%%%%%%%%%%%%%%%%%%%%%%%%%%%%%%%%%%%%%%%%%%%%%%%%
% A Beamer template for University of Wollongong     %
% Based on THU beamer theme                          %
% Author: Qiuyu Lu                                   %
% Date: July 2024                                    %
% LPPL Licensed.                                     %
%%%%%%%%%%%%%%%%%%%%%%%%%%%%%%%%%%%%%%%%%%%%%%%%%%%%%%
% Customized for Sharif University of Technology     %
%%%%%%%%%%%%%%%%%%%%%%%%%%%%%%%%%%%%%%%%%%%%%%%%%%%%%%


\documentclass[serif, aspectratio=169]{beamer}
%\documentclass[serif]{beamer}  % for 4:3 ratio
\usepackage[T1]{fontenc} 
\usepackage{fourier} % see "http://faq.ktug.org/wiki/uploads/MathFonts.pdf" for other options
\usepackage{hyperref}
\usepackage{latexsym,amsmath,xcolor,multicol,booktabs,calligra}
\usepackage{graphicx,pstricks,listings,stackengine}
\usepackage{lipsum}
\usepackage{tikz}
\usetikzlibrary{shapes.geometric} % Add this to include ellipse shapes
\usepackage{amsmath,amsthm}
\usepackage{pgfplots}  % For plots
\usepackage{amsmath}   % For equations
\usepackage{array}     % For tables
\pgfplotsset{compat=1.16}

% \usetikzlibrary{positioning}


%\newtheorem{theorem}{Theorem}
%\newtheorem{definition}{Definition}
%\newtheorem{remark}{Remark}


\newcommand{\relu}{\text{ReLU}}
\newcommand{\E}{\mathbb{E}}


\author{Hooman Zolfaghari - Abdollah Zohrabi - Amirreza Velae}
\title{Neural Tangent Kernel}
\subtitle{High-Dimensional Probability Analysis}
\institute{
    Sharif University of Technology
}
%\date{\small \today}
% \usepackage{UoWstyle}
\usepackage{SUTstyle}

% defs
\def\cmd#1{\texttt{\color{red}\footnotesize $\backslash$#1}}
\def\env#1{\texttt{\color{blue}\footnotesize #1}}
\definecolor{deepblue}{rgb}{0,0,0.5}
\definecolor{deepred}{RGB}{153,0,0}
\definecolor{deepgreen}{rgb}{0,0.5,0}
\definecolor{halfgray}{gray}{0.55}

\lstset{
    basicstyle=\ttfamily\small,
    keywordstyle=\bfseries\color{deepblue},
    emphstyle=\ttfamily\color{deepred},    % Custom highlighting style
    stringstyle=\color{deepgreen},
    numbers=left,
    numberstyle=\small\color{halfgray},
    rulesepcolor=\color{red!20!green!20!blue!20},
    frame=shadowbox,
}

\begin{document}

\begin{frame}
    \titlepage
    \vspace*{-0.6cm}
    \begin{figure}[htpb]
        \begin{center}
            \includegraphics[keepaspectratio, scale=0.25]{pic/sharif-main-logo.png}
        \end{center}
    \end{figure}
\end{frame}

\begin{frame}    
\tableofcontents[sectionstyle=show,
subsectionstyle=show/shaded/hide,
subsubsectionstyle=show/shaded/hide]
\end{frame}

% ============================ Introduction ============================ 
\section{Introduction}

%-------------------------------------------------
\begin{frame}{Key Idea}
	\begin{itemize}
		\item \textbf{Dynamics} of training infinitely wide NNs \(\approx\)
		\item \textbf{convex optimization} in RKHS
		
		\vspace{0.5em}
		\item For \(f_\theta: \mathbb{R}^d \to \mathbb{R}\) (NN), GD training induces:
		\[
		\underbrace{\Theta(\theta)_{i,j}}_{\text{NTK}} := \nabla_\theta f_\theta(x_i) \nabla_\theta f_\theta(x_j)^\top \in \mathbb{R}^{n \times n}
		\]
		\small
		\textbf{Crucial Property}: \(\Theta(\theta^{(0)}) \to \Theta^{\infty}\) as width \(\to \infty\)
		
	\end{itemize}
	
	
\end{frame}

%-------------------------------------------------
\begin{frame}{Modern Perspective}
	\begin{itemize}
		\item \textbf{Feature Learning Gap}: NTK regime \(\neq\) real NNs (finite-width trains via \(\nabla\Theta \neq 0\))
		\item \textbf{Deep vs Shallow}: Depth induces \textbf{spectral bias} (eigenvalue decay of \(\Theta^\infty\))
	\end{itemize}
	
	\[
	\partial_t f_t = -\Theta^\infty(f_t - y) \quad (\text{grad flow ODE})
	\]
	
	\vspace{1em}
	\begin{columns}
		\begin{column}{0.5\textwidth}
			\textbf{Good News}:
			\begin{itemize}
				\item Global convergence
				\item Linear rate for \(\lambda_{\min}(\Theta^\infty) > 0\)
			\end{itemize}
		\end{column}
		\begin{column}{0.5\textwidth}
			\textbf{Bad News}:
			\begin{itemize}
				\item No feature learning
				\item Fails for transformers/attention
			\end{itemize}
		\end{column}
	\end{columns}
\end{frame}

\section{Analysis of Convergence and Generalization}

\begin{frame}
	\begin{itemize}
		\item Research focused on more practical assumptions.
		\item Two-layer ReLU network \(f_{\mathbf{W},\mathbf{a}}(\mathbf{x}) = \frac{1}{\sqrt{m}}\sum_{r=1}^{m} a_r \relu(\mathbf{w}_r^\top \mathbf{x})
		\)
		\item Least Squares Regression \(C(W):=\frac{1}{2}\sum_{i=1}^{n}(y_i - f_{\mathbf{W},\mathbf{a}}(\mathbf{x}_i))\)
		\item Resulting NTK gram matrix:
		\begin{align}
		\mathbf{H}_{ij}^\infty &= \E{\mathbf{w} \sim \mathcal{N}(0,\mathbf{I})}\left[ \mathbf{x}_i^\top \mathbf{x}_j\mathbb{I}\left\{\mathbf{w}^\top \mathbf{x}_i \ge 0, \mathbf{w}^\top \mathbf{x}_j \ge 0\right\}\right] \\
		&= \frac{\mathbf{x}_i^\top\mathbf{x}_j\left(\pi - \arccos(\mathbf{x}_i^\top\mathbf{x}_j)\right)}{2\pi}, \quad \forall i, j\in[n].
		\end{align}
		\item  If \(H^\infty\) is positive
		definite \(\lambda_0 := \lambda_{\min}(H^{\infty})>0\) , GD converges to \(0\) training loss if \(m\) is sufficiently large \(\Omega(\frac{n^6}{\lambda_0^4})\).
	\end{itemize}
\end{frame}



%----------------------------------------------------
%  Eigen-decomposition setup

\begin{frame}
\begin{itemize}

		\item Eigen-decomposition $H^\infty\;=\; \sum_{i=1}^n \lambda_i\,v_i\,v_i^\top,$.

%----------------------------------------------------
%  Theorem 4.1

\item Suppose $\lambda_0 = \lambda_{\min}(H^\infty) \;>\; 0$, 
	\(
	\kappa \;=\; O\!\biggl(\frac{\varepsilon_0\,\delta}{\sqrt{n}}\biggr),
	\quad
	m \;=\; \Omega\!\Bigl(\frac{n^7}{\lambda_0^4\,\kappa^2\,\delta^4\,\varepsilon^2}\Bigr),
	\quad
	\eta \;=\; O\!\Bigl(\frac{\lambda_0}{n^2}\Bigr).
	\)
	Then with probability at least $1 - \delta$ over the random initialization, 
	for all $k = 0,1,2,\dots$ we have
	\[
	\|y \;-\; u(k)\|_2 
	\;=\;
	\sqrt{\;\sum_{i=1}^n \Bigl(1 \;-\; \eta\,\lambda_i\Bigr)^{2k}\,\bigl(v_i^\top y\bigr)^2}
	\;\;\pm\;\;\varepsilon.
	\tag{8}
	\]


%----------------------------------------------------
%  Definition 5.1

\item 
	A distribution $D$ over $\mathbb{R}^d \times \mathbb{R}$ is called 
	$(\lambda_0,\delta,n)$\emph{-non-degenerate} 
	if for $n$ i.i.d.\ samples $\{(x_i,y_i)\}_{i=1}^n$ from $D$, 
	with probability at least $1 - \delta$ we have
	\[
	\lambda_{\min}\bigl(H^\infty\bigr) \;\;\ge\;\; \lambda_0 \;>\; 0.
	\]

\end{itemize}

\end{frame}


%----------------------------------------------------
%  Remark 5.1
%
%\begin{remark}[Remark 5.1]
%	\label{rem:5.1}
%	Note that as long as no two $x_i$ and $x_j$ are parallel to each other, 
%	we have $\lambda_{\min}\bigl(H^\infty\bigr) > 0$. 
%	For most real-world distributions, any two training inputs are not parallel.
%\end{remark}

%----------------------------------------------------
%  Theorem 5.1

\begin{frame}
	
	\label{thm:5.1}
	Fix a failure probability $\delta \in (0,1)$. 
	Suppose our data $S = \{(x_i,y_i)\}_{i=1}^n$ are i.i.d.\ samples 
	from a $(\lambda_0,\delta/3,n)$-non-degenerate distribution $D$, 
	and let 
	\[
	\kappa \;=\; O\!\Bigl(\tfrac{\lambda_0\,\delta}{n}\Bigr),
	\quad
	m \;\ge\; \kappa^{-2}\,\mathrm{poly}\bigl(n,\,\lambda_0^{-1},\,\delta^{-1}\bigr).
	\]
	Consider any loss function 
	$\ell:\mathbb{R}\times\mathbb{R}\to[0,1]$ 
	that is 1-Lipschitz in its first argument and satisfies $\ell(y,y)=0$. 
	Then with probability at least $1 - \delta$ 
	(over the random initialization \emph{and} the training samples), 
	the two-layer neural network $f_{\mathbf{W}(k),a}$ 
	trained by gradient descent for 
	\[
	k \;\ge\; \Omega\!\Bigl(\tfrac{1}{\eta\,\lambda_0}\,\log\tfrac{n}{\delta}\Bigr)
	\quad\text{iterations}
	\]
	has population loss
	\[
	L_D\bigl(f_{\mathbf{W}(k),a}\bigr)
	\;=\;
	\mathbb{E}_{(x,y)\sim D}\bigl[\ell\bigl(f_{\mathbf{W}(k),a}(x),\,y\bigr)\bigr]
	\;\le\;
	\sqrt{\frac{2\,y^\top \bigl(H^\infty\bigr)^{-1}\,y}{n}}
	\;+\;
	O\!\Bigl(\sqrt{\tfrac{\log\!\bigl(\tfrac{n}{\lambda_0\,\delta}\bigr)}{n}}\Bigr).
	\tag{9}
	\]

\end{frame}

\section{Closer Look \& Motivation}

\begin{frame}
	\begin{itemize}
		\item We can see that the bound depends on:
		\begin{itemize}
		
			 \item Distribution \((x,y) \sim \mathcal{D}\) such that,
			\item \(\mathbf{y}^\top H^\infty \mathbf{y} \leq \|(H^\infty)^{-1}\| \|\mathbf{y}\|_2 = \lambda_{\min}(H^{\infty}) \|\mathbf{y}\|_2 \)
		\end{itemize}  
		
		
		\item To be able to \text{provably} learn (e.g. PAC-Learning), the bound must converge to \(0\) as \(n \to \infty\).
		
		\item The paper mentions the case of \(y=g(x)\) for some function \(g\).
		
		\item So we focus on bounding \(\lambda_{\min}(H^{\infty}) \text{ and } \|\mathbf{y}\|_2\)
		
	\end{itemize}

\end{frame}


\section{Bounds on Minimum Eigenvalue}


\begin{frame}
\begin{itemize}
	\item Data scaling assumption.
		The data distribution $P_X$ satisfies the following properties:
		\begin{enumerate}
			\item $\displaystyle \int \|x\|_{2}\, dP_X(x) \;=\; \Theta(\sqrt{d}).$
			\item $\displaystyle \int \|x\|_{2}^{2}\, dP_X(x) \;=\; \Theta(d).$
			\item $\displaystyle \int \bigl\|x \;-\; \int x' \, dP_X(x')\bigr\|_{2}^{2}\, dP_X(x) 
			\;=\; \Omega(d).$
		\end{enumerate}
		\item These are just scaling conditions on the data vector $x$ or its centered 
		counterpart $x - \mathbb{E}x$. We remark that the data can have any scaling, 
		but in this paper we fix it to be of order $d$ for convenience. We further 
		assume the following condition on the data distribution.
	

\end{itemize}
\end{frame}

\begin{frame}
	\begin{itemize}
			\item Lipschitz concentration assumption.
		The data distribution $P_X$ satisfies the Lipschitz concentration property. 
		Namely, for every Lipschitz continuous function 
		$f : \mathbb{R}^d \to \mathbb{R}$, there exists an absolute constant $c > 0$ 
		such that, for all $t > 0$,
		\[
		\mathbb{P}\Bigl(\bigl|f(x) - \int f(x')\,dP_X(x')\bigr| \;>\; t\Bigr)
		\;\;\le\;\; 2\,e^{-\,c\,t^2\,/\,\|f\|_{\mathrm{Lip}}^2}.
		\]
		\item In general, this assumption covers the whole family of distributions that satisfy 
		the log-Sobolev inequality with a dimension-independent constant (or distributions
		with log-concave densities).
	\end{itemize}
\end{frame}


\begin{frame}
	\begin{theorem}[Smallest eigenvalue of limiting NTK]
		\label{thm:smallest-eig-NTK}
		Let $\{x_i\}_{i=1}^N$ be a set of i.i.d.\ data points from $P_X$, where $P_X$ has 
		zero mean and satisfies Assumptions~2.1 and~2.2. Let $K^{(L)}$ be the limiting 
		NTK recursively defined in (9). Then, for any even integer constant $r \ge 2$, 
		we have with probability at least 
		\[
		1 - N\,e^{-\Omega(d)} \;-\; N^2\,e^{-\Omega\bigl(d\,N^{-2/(\,r - 0.5\,)}\bigr)}
		\]
		that
		\[
		\mathrm{LO}(d) 
		\;\;\ge\;\;
		\lambda_{\min}\bigl(K^{(L)}\bigr) 
		\;\;\ge\;\;
		\mu_{r}(\sigma)^{2}\;\Omega(d),
		\]
		where $\mu_{r}(\sigma)$ is the $r$-th Hermite coefficient of the ReLU function 
		given by (8).
	\end{theorem}
\end{frame}


\section{Our Results and Observations}

\begin{frame}
\begin{itemize}
	  \item Most papers and our main reference assume \(\|x\|=1\) and \(|y|\leq 1\), for simplicity.
	
	\item Therefore, the diagonal \(H^\infty_{ii} = \frac{1}{1}\).
	\item Now denote \(\rho = \max_{i,j\neq i} |x_i^\top x_j| \). since \(f(x)=\frac{x(\pi-\arccos(x))}{2\pi}\), is as below, we get:
	\[
	H^\infty_{i,j\neq i} \leq \frac{\rho(\pi-\arccos(\rho))}{2\pi} \leq \frac{1}{2}
	\]
	\item We can use the \textbf{Gershgorin circle} theorem and get
	\[
	| \lambda_{\min}(H^\infty) - \frac{1}{2} | \leq (n-1)\frac{\rho(\pi-\arccos(\rho))}{2\pi}
	\]
	\item Thus, the bound depends on the maximum "correlation" between two distinct i.i.d \(x_i\) in a sample of size \(n\). This will depend on the distribution of \(x\) on \(S^{d-1}\).
\end{itemize}
\end{frame}


\begin{frame}
	\begin{itemize}

	\item So the bound is good when \(x_i\) are almost orthogonal with high probability. We can show the following examples:
	
	\item \(x_i \sim \text{Unif}(S^{d-1})\):
	\[
	\max_{1 \le i < j \le n}
	\Bigl|\langle X_i,\;X_j\rangle\Bigr|
	\;\;\le\;
	C\,
	\sqrt{\frac{\log\!\bigl(\tfrac{n}{\delta}\bigr)}{d}}
	\,.
	\]
	
	\item \(x_i\) are isotropic, mean-zero, sub-Gaussian vectors:
	\[
	\max_{1 \le i < j \le n}
	\bigl|\langle X_i,\;X_j\rangle\bigr|
	\;\;\le\;
	C\,
	\sqrt{\frac{\log\!\bigl(\tfrac{n}{\delta}\bigr)}{d}}
	\;\;\max_{1\le i\le n}\|X_i\|\ .
	\]
	\end{itemize}
	
\end{frame}


\begin{frame}
	 Bernstein Result
\end{frame}




%%-------------------------------------------------
%\begin{frame}{Advanced Implications}
%	\begin{itemize}
%		\item \textbf{Phase Transitions}: \(\exists\) critical widths where NTK/feature learning dominates
%		\item \textbf{Quantum Analog}: NTK \(\leftrightarrow\) quantum metric tensor (geometry of NN L2 space)
%		\item \textbf{Diffusion Approx}: SGD noise \(\propto \Theta_\infty^{-1}\) (implicit regularization)
%	\end{itemize}
%	
%	\vspace{1em}
%	\textbf{Open Problems}:
%	\begin{itemize}
%		\item NTK for \textbf{dynamical} architectures (RNNs, SSMs)
%		\item \textbf{Non-Gaussian} initialization limits
%	\end{itemize}
%\end{frame}

%-------------------------------------------------


\section{References}




%\begin{frame}{Contributions}
%\textbf{These slides are authored by:}
%    \begin{itemize}
%        \item Hooman Zolfaghari
%    \end{itemize}
%    
%\end{frame}


\begin{frame}[allowframebreaks]
   \bibliographystyle{ieeetr} % Place the style before bibliography
   \bibliography{ref.bib} % Point to your .bib file
   \nocite{*} % Include all references even if not cited
\end{frame}


\end{document}